%%FILE I-a
%%APPEND title="Equivalence measured by constant coefficients"
How much less than the homotopy type a coefficient theory detects is computed for sets, truncated types, complexes over a ring and spectra in the commentary on Chapter~X, ms.~5, together with the finer class defined by locally constant coefficients. That finer class is a fundamental localizer for every derivator (commentary at ms.~100 of this chapter). The strict $\mathbb D$-equivalences of ms.~65 have the closure properties of the intermediate, 1-complete class: composition and the rule that $f$ and $gf$ in the class put $g$ in it, without two-out-of-three.

%%FILE VI-a
%%APPEND title="The locally constant comparison is now a theorem"
Two facts hold for every derivator on all small categories, with no hypothesis on the localizer: the locally constant objects of $\mathbb D(S)$ form a full subcategory whose inclusion has both adjoints, compatibly with change of the diagram category; and a functor $A\to B$ over $S$ that induces equivalences on locally constant objects over every slice $S/s$ induces one on $\mathrm{LC}(B)\to\mathrm{LC}(A)$. The first is the theorem of Cisinski cited above; the second is the relative Theorem~A for the localizer of locally constant coefficients, proved in the commentary on Chapter~I, ms.~100. They concern the localizer $W^{\mathrm{lc}}_{\mathbb D}$ attached to the coefficients, which for $\mathsf{HOT}$ is $W_\infty$; they do not extend the identifications of this chapter to an arbitrary $W$.

%%FILE XIIc-a
%%APPEND title="Degree two and the localizer correction"
The three classes have a meaning in terms of coefficients. For the derivator of $n$-truncated homotopy types the constant-coefficient equivalences are the maps satisfying $b_n$; the functors whose restriction is fully faithful on locally constant coefficients are those satisfying $a_n$; and those whose restriction is an equivalence there satisfy $b_{n+1}$ (table in the commentary on Chapter~X, ms.~5). So $b_{n+1}\Rightarrow a_n\Rightarrow b_n$ is an instance of the general chain of the commentary on Chapter~I, ms.~100, and the failure of two-out-of-three for $a_n$ is the failure for the 1-complete class there: only the two outer classes are localizers.

%%FILE XI-XIIa
%%APPEND title="The universal meaning of a Hot-fibration"
For the derivator of spaces the classes of equivalences of Chapter~I all coincide with $W_\infty$, and the fibration conditions of that chapter (the universal cohomological one, the homological one, the test over maps between weakly contractible bases) are then expected to coincide with universal preservation of $W_\infty$ under base change; for another localizer the test over contractible bases is a different notion (commentary on Chapter~I, ms.~118 and~122). The coincidence for spaces is recorded here as stated; it was not verified for this commentary.

%%FILE XVI
%%APPEND title="A reverse-excision counterexample in every finite degree"
The same failure occurs for every localizer $W^{\mathrm{lc}}_{\mathbb D}$ of locally constant coefficients of a nontrivial derivator $\mathbb D$ (commentary on Chapter~I, ms.~100): none satisfies (6b) or (6c), whereas all satisfy (6a). When $\mathbb D(e)$ is not a preorder the test is the double cone on $BG$ for a nontrivial acyclic group $G$, which is contractible while $BG\to e$ is not in the class; when it is a preorder, the fold of two points, whose double cone is a circle. The verification is not reproduced here.

%%APPEND title="No homology localizer has enough intrinsic fibrations"
\textbf{A further example: rational local coefficients.}\ours
Let $V$ be the class of functors inducing a bijection on $\pi_0$, isomorphisms on $\pi_1$ and isomorphisms on $\pi_j\otimes\mathbf Q$ for $j\geq2$; it is the localizer of locally constant coefficients in complexes of rational vector spaces (commentary on Chapter~X, ms.~5), hence a fundamental localizer satisfying (3 bis). It does not satisfy Loc~(7F).

\textit{Proof.} A $V$-aspherical category $A$ is connected and simply connected with torsion homotopy groups, so its reduced integral homology is torsion and $H^1(A;\mathbf Z)=H^2(A;\mathbf Z)=0$. Hence the constants $K(\mathbf Z,2)\to\operatorname{Map}(BA,K(\mathbf Z,2))$ form a weak equivalence, and $A\to e$ belongs to the fundamental localizer $W'$ of the functors $f$ for which $\operatorname{Map}(Bf,K(\mathbf Z,2))$ is a weak equivalence. If $V$ satisfied (7F), the second paragraph of the proof of Proposition~2 at ms.~27, which uses only (3 bis) and (7F), would give $V\subset W'$. But multiplication by a prime $p$ on $K(\mathbf Z,2)$ is in $V$, and it acts on $[K(\mathbf Z,2),K(\mathbf Z,2)]=\mathbf Z$ by multiplication by $p$.\hfill$\square$

More generally, for a nonzero ring $k$ the localizer of locally constant coefficients in complexes of $k$-modules is stated to satisfy (7F) exactly when no prime is invertible in $k$, that is, exactly when it is $\mathbf W_\infty$; only the case $k=\mathbf Q$ is proved here.
